\section{Introduction}
\label{sec:intro}
Surface charging of airless bodies plays a crucial role in shaping the near-surface plasma environment.
Photoelectron emission dominates the current balance on the dayside, whereas solar wind electrons dominate on the nightside \cite{Whipple1981,Nitter1998}.
However, the role of small-scale topography remains poorly understood.
Previous particle-in-cell (PIC) studies have shown that crater walls can charge to \wallPot{}~V relative to the ambient plasma \cite{Zimmerman2016}.
In this paper, we delve into how crater size controls the surface potential.
